\documentclass[11pt, a4paper]{article}
\usepackage{amsmath}
\usepackage{amssymb}
\usepackage{graphicx}
\usepackage{hyperref}
\usepackage{listings}
\usepackage[margin=1in]{geometry}
\usepackage{fancyhdr}
\usepackage{booktabs}
\usepackage{tabularx}

\lstset{
    basicstyle=\ttfamily\small,
    keywordstyle=\color{blue},
    commentstyle=\color{gray},
    breaklines=true,
    language=Python
}

\title{ECG Simulator Documentation\\Part 6: System Architecture \& Design}
\author{Module Organization, Class Hierarchy, Extension Points}
\date{June 2026}

\pagestyle{fancy}
\fancyhf{}
\rhead{Part 6: Architecture}
\lhead{ECG Simulator}
\cfoot{\thepage}

\begin{document}

\maketitle
\tableofcontents
\newpage

\section{Three-Layer Architecture}

\subsection{Overview}
The ECG simulator is organized into three distinct layers:

\begin{enumerate}
    \item \textbf{Physics Layer}: Cardiac physiology (cell models, propagation, dipoles)
    \item \textbf{Simulation Layer}: Algorithm selection (Phase 1/2 vs Phase 3), beat management
    \item \textbf{GUI Layer}: PyQt6 interface, real-time visualization, parameter control
\end{enumerate}

\subsection{Layer Interaction}

\begin{verbatim}
GUI Layer (main_window.py, parameter_panel.py)
    |
    | controls
    |
Simulation Layer (graph_engine.py, cable_engine.py)
    |
    | uses
    |
Physics Layer (cell_models/, tissue/, ecg/)
\end{verbatim}

\subsection{Data Flow}
\begin{enumerate}
    \item User adjusts HR or parameters in GUI
    \item GUI updates ParameterModel (single source of truth)
    \item SimulationEngine queries ParameterModel for each beat
    \item Engine executes propagation algorithm
    \item Forward model computes 12-lead ECG voltages
    \item GUI displays voltages on pyqtgraph widget
\end{enumerate}

\section{Physics Layer (Core Modules)}

\subsection{Module Organization}

\begin{tabularx}{\textwidth}{|l|X|c|}
\hline
\textbf{Module} & \textbf{Purpose} & \textbf{Lines} \\
\hline
cardiac\_sim/core/cell\_models/fitzhugh\_nagumo.py & FHN cell model ODEs & 150 \\
cardiac\_sim/core/tissue/cable\_1d.py & 1-D monodomain integration & 500 \\
cardiac\_sim/core/tissue/graph.py & DAG + BFS conduction & 350 \\
cardiac\_sim/core/ecg/lead\_field.py & 12-lead lead vectors & 150 \\
\hline
\end{tabularx}

\subsection{Class Hierarchy: Cell Models}

\begin{verbatim}
BaseCellModel (abstract)
    ├── FitzHughNagumoCell
    └── [Future: HodgkinHuxley, Courtemanche, tenTusscher]

Key Methods:
  - get_initial_state()
  - compute_derivatives(t, state)
  - stimulate(t_start, duration, amplitude)
  - v_to_mv(v_norm, v_rest, v_peak)
\end{verbatim}

\subsubsection{FitzHughNagumoCell Interface}

\begin{lstlisting}
class FitzHughNagumoCell(BaseCellModel):
    def get_initial_state(self) -> np.ndarray:
        # Returns [V_REST, W_REST]
        return np.array([-1.200, -0.625])
    
    def compute_derivatives(self, t, state):
        # Solves dimensionless FHN ODEs
        # Scales to physical time via TAU_S
        v, w = state
        dv = v - v**3/3 - w + I_ext
        dw = self.eps * (v + self.a - self.b * w)
        return np.array([dv, dw])
    
    def v_to_mv(self, v_norm, v_rest, v_peak):
        # Maps [-1.2, 2.0] → [-85, +30] mV
        return v_rest + (v_peak - v_rest) * \
               (v_norm - self.V_REST) / (self.V_PEAK - self.V_REST)
\end{lstlisting}

\subsection{Cable Model (1-D Tissue)}

\begin{lstlisting}
class Cable1D:
    def __init__(self, params):
        # Initialize V, W state arrays (18 cells)
        self.V = np.ones(18) * self.V_REST
        self.W = np.ones(18) * self.W_REST
        self._check_cfl()
    
    def new_beat(self, sa_fire_time, params):
        # Reset state, mark SA activated
        self.V.fill(self.V_REST)
        self.W.fill(self.W_REST)
        self._active_beat_start = sa_fire_time
    
    def advance(self, dt):
        # Integrate dt seconds via n_micro steps
        # Returns dict of newly activated nodes
        for _ in range(n_micro):
            self._step_atrial(DT_INT)
            self._step_ventricular(DT_INT)
        return newly_activated_dict
\end{lstlisting}

\section{Simulation Layer}

\subsection{SimulationEngine Interface}

\begin{verbatim}
SimulationEngine (abstract)
    ├── GraphEngine (Phase 1/2: BFS-based)
    └── CableEngine (Phase 3: FHN cable)

Key Methods:
  - generate_beat(hr, params) -> dict[node -> t_act]
  - get_activation_times() -> dict
  - advance_time(dt)
\end{verbatim}

\subsection{GraphEngine (Phase 1/2)}

\begin{verbatim}
BFS Algorithm:
1. Start queue with SA node (t=0)
2. For each node:
   a. Compute conduction time to each neighbor
   b. Check if neighbor is refractory
   c. If NOT refractory: add to queue with new activation time
   d. If refractory: skip
3. Return dict[node -> activation_time] for entire beat

Refractory Check:
  t_candidate = t_current + conduction_time
  if t_candidate < last_activation[node] + ERP[node]:
      # Skip (still refractory)
  else:
      # Conduct
\end{verbatim}

\subsubsection{Key Implementation Details}

\begin{lstlisting}
def run_bfs(self, sa_fire_time, params):
    activation_times = {}
    queue = deque([(SA_NODE, sa_fire_time)])
    
    while queue:
        current_node, t_curr = queue.popleft()
        activation_times[current_node] = t_curr
        
        # Find refractory status
        last_act = self._last_activation.get(current_node, -inf)
        if t_curr < last_act + ERP[current_node]:
            continue  # Skip if still refractory
        
        # Process neighbors
        for neighbor, cond_time in self.adjacency[current_node]:
            t_next = t_curr + cond_time
            queue.append((neighbor, t_next))
    
    return activation_times
\end{lstlisting}

\subsection{CableEngine (Phase 3)}

\begin{verbatim}
Phase 3: FHN Cable Integration
1. Initialize 18-cell cable (8 atrial + 10 ventricular)
2. Stimulate SA node at t=sa_fire_time
3. Advance cable.advance(dt) every 33 ms
4. Cable returns dict of newly activated nodes
5. Accumulate activation times incrementally
6. Compute dipole moments from activation_times
7. Project to 12 leads

Architecture:
  - _cable: Cable1D object (manages state V, W)
  - _active_beats: List of BeatRecord objects
  - _current_beat_nodes: Dict tracking activated nodes per beat
  - _nodes: 18 ECGNode objects (anatomical positions)
  - _forward_model: LeadField for projection
\end{verbatim}

\section{GUI Layer}

\subsection{PyQt6 Widget Hierarchy}

\begin{verbatim}
MainWindow (QMainWindow)
    ├── ControlPanel (QWidget)
    │   ├── HR Spinner (QSpinBox)
    │   ├── Engine Selector (QComboBox)
    │   ├── Control Buttons (QPushButton)
    │   └── Display Label (QLabel)
    ├── ECGDisplay (PlotWidget from pyqtgraph)
    │   └── 12 curves (one per lead)
    └── ParameterPanel (QScrollArea)
        ├── SANodeGroup (QGroupBox)
        ├── AVNodeGroup (QGroupBox)
        ├── ... (12 groups total)
        └── ApplyButton (QPushButton)
\end{verbatim}

\subsection{Signal/Slot Architecture}

\begin{lstlisting}
# User changes HR spinbox
spinbox.valueChanged.connect(self.on_hr_changed)

def on_hr_changed(self, new_hr):
    self.params.sa.cycle_length = 60 / new_hr
    # Changes take effect on next beat

# Timer updates ECG display
self.timer.timeout.connect(self.update_display)

def update_display(self):
    # Advance simulation
    self.engine.advance(DT_FRAME)
    
    # Get ECG voltages
    ecg_data = self.engine.get_ecg_data()
    
    # Update all 12 curves
    for i, lead in enumerate(self.leads):
        self.plot.getPlotItem().curves[i].setData(ecg_data[i])
\end{lstlisting}

\section{Plugin System}

\subsection{Plugin Architecture}

\begin{verbatim}
PluginBase (abstract)
    ├── StaticPlugin (cached, whole-beat effect)
    └── StatefulPlugin (per-beat, temporary override)

Stacking Mechanism:
1. Phase 1/2 BFS algorithm runs
2. For each plugin in stack:
   - Modify activation_times dict or conductances
   - Can disable regions (set conductance → 0)
   - Can prolong refractories
3. Result: Composite arrhythmia
\end{verbatim}

\subsection{Available Plugins}

\begin{tabularx}{\textwidth}{|l|X|}
\hline
\textbf{Plugin} & \textbf{Mechanism} \\
\hline
AVBlockI & Prolong av\_delay \\
AVBlockII\_Mobitz1 & Reduce av\_conductance gradually \\
AVBlockII\_Mobitz2 & Reduce his\_conductance abruptly \\
AVBlockIII & Set av\_conductance → 0 \\
LBBB & Set lbb\_conductance → 0 \\
RBBB & Set rbb\_conductance → 0 \\
LAHB & Set lahb\_conductance → 0 \\
LPHB & Set lphb\_conductance → 0 \\
AF & Insert fibrillatory waves into atrial rate \\
PVC & Inject ectopic stimulus at coupling interval \\
Sinus Bradycardia & Increase sa cycle\_length \\
\hline
\end{tabularx}

\section{Data Flow Example: Normal Beat}

\begin{enumerate}
    \item User sets HR = 70 bpm
    \item GUI calculates cycle\_length = 60/70 ≈ 0.857 s
    \item Timer tick at t=0: SA fires
    \item BFS/Cable propagates: SA → RA (124 ms) → LA (191 ms) → AV (307 ms) → His (435 ms) → LV (645 ms)
    \item Dipole computation for each region
    \item Lead field projection: 18-D position vector → 12-lead voltages
    \item Display updates all 12 curves
    \item Next cycle_length: SA fires again
\end{enumerate}

\section{Performance Characteristics}

\subsection{Computational Complexity}

\begin{tabularx}{\textwidth}{|l|c|c|}
\hline
\textbf{Engine} & \textbf{BFS Complexity} & \textbf{Actual Rate} \\
\hline
Phase 1/2 & O(V + E) = O(18 + edges) & >1000 Hz \\
Phase 3 & O(N_cells × N\_steps) & 2.2× real-time \\
\hline
\end{tabularx}

\subsection{Memory Usage}

\begin{itemize}
    \item State arrays: ~18 cells × 2 variables (V, W) × 8 bytes = 288 bytes
    \item ECG buffer: ~10 seconds × 30 FPS × 12 leads × 8 bytes = 28.8 KB
    \item Activation times: ~18 nodes × 100 beats × 8 bytes = 14.4 KB
    \item Total: <100 MB for typical 1-hour simulation
\end{itemize}

\section{Extension Points}

\subsection{Adding a New Cell Model}

\begin{enumerate}
    \item Create subclass of BaseCellModel
    \item Implement: get\_initial\_state(), compute\_derivatives()
    \item Add to cable\_1d.py imports
    \item Update engine selector to include new model
\end{enumerate}

\subsection{Adding a New Plugin}

\begin{enumerate}
    \item Create class inheriting PluginBase
    \item Implement: apply() method (modifies activation times or conductances)
    \item Register in plugin manager
    \item Add to GUI parameter panel
\end{enumerate}

\subsection{Adding a New Lead Vector}

\begin{enumerate}
    \item Define unit vector (X, Y, Z) in cardiac coordinate system
    \item Add to LEAD\_MATRIX in lead\_field.py
    \item Update lead names list
    \item GUI automatically includes in display
\end{enumerate}

\end{document}
